% !TeX program = xelatex
\documentclass[12pt]{article}
\usepackage{fontspec}
\setmainfont{Liberation Serif}   % oder Times New Roman, Arial, etc.
\usepackage[ngerman]{babel}      % deutsche Silbentrennung und Lokalisierung

\usepackage{amsmath,amssymb,amsthm,mathtools}
\usepackage{geometry}
\usepackage{booktabs}
\usepackage{hyperref}
\usepackage{url}
\usepackage{tabularx}
\usepackage{array}
\usepackage{makecell}
\usepackage{ragged2e}
\usepackage{bm}
\usepackage{graphicx}
\usepackage{caption}
\usepackage{subcaption}
\usepackage{float}

\newtheorem{theorem}{Satz}
\newtheorem{lemma}{Lemma}
\newtheorem{definition}{Definition}
\newtheorem{corollary}{Korollar}
\theoremstyle{remark}
\newtheorem{remark}{Bemerkung}

\newcommand{\Keff}{K_{\mathrm{eff}}}
\newcommand{\eps}{\varepsilon}
\newcommand{\df}{d_{\!f}}
\newcommand{\kappac}{\kappa}
\newcommand{\betaf}{\beta}
\newcommand{\Tmin}{T_{\min}}
\newcommand{\Dprior}{D_{\mathrm{prior}}}

\hypersetup{colorlinks=true,linkcolor=blue,citecolor=blue,urlcolor=blue}

\begin{document}
	
	\begin{titlepage}
		\begin{center}
			\vspace*{0.5cm}
			\Huge\textbf{Anhang O. Universelle Skalierung kritischer Gruppengrößen: \\ Empirische Evidenz und Verbindung zum Axiom der technologischen Schwelle (YUCT)}
			
			\vspace{0.3cm}
			\small\textit{Eine systematische Sammlung von Daten über kritische Gruppengrößen in biologischen und sozialen Systemen, interpretiert durch das Axiom der technologischen Wörterschwelle (Anhang O) und den universellen Exponenten $\beta = 0,67$ (Anhang L).}
			
			\vspace{0.3cm}
			\small\textbf{Alexey V. Yakushev} \\
			YUCT-Theorie: \url{https://yuct.org/}\\
			YPSDC-Protokoll: \url{https://ypsdc.com/}\\
			
			\vspace{2cm}
			\textbf DOI: \url{https://doi.org/10.5281/zenodo.18444598}\\
			
			\vspace{1cm}
			
			\vspace{0cm}
			\large 
			Eine Kurzversion dieser Arbeit wurde zuvor veröffentlicht und ist verfügbar unter\\ \url{https://doi.org/10.5281/zenodo.18362308}\\
			\vspace{1cm}
			März 2026 \\ \large Version 2.0.
			
			\vspace{0.5cm}
			\begin{minipage}{0.95\textwidth}
				\begin{abstract}
					\noindent
					Dieser Anhang fasst empirische Daten zu den kritischen Größen zusammen, bei denen verschiedene kollektive Systeme einen Phasenübergang durchlaufen – wie Schwärmen, Aufspaltung oder strukturelle Reorganisation. Untersucht werden Honigbienenvölker, Militäreinheiten, Fischschwärme, Delfingruppen, Nutztierpopulationen und Primatengruppen. Unter Verwendung des universellen Fehlerskalierungsgesetzes $\varepsilon \propto N^{\beta}$ mit $\beta = 0,67$ (Anhang L) und des Axioms der technologischen Wörterschwelle (Anhang O) leiten wir eine quantitative Beziehung zwischen der kritischen Größe $N_{\text{crit}}$, bei der ein Übergang stattfindet, und der optimalen Gruppengröße $N_{\text{opt}}$ ab. Die beobachteten Verhältnisse $N_{\text{crit}}/N_{\text{opt}}$ für verschiedene Systeme gruppieren sich um Werte, die mit $\beta = 0,67$ konsistent sind, und liefern so eine unabhängige Stütze für die Universalität dieses Exponenten sowie für die Idee, dass das Überschreiten einer technologischen (oder organisatorischen) Schwelle einem Phasenübergang der Koordinationseffizienz entspricht. Wir schlagen außerdem spezifische, überprüfbare Vorhersagen für zukünftige Experimente und Feldstudien vor, die den YUCT-Rahmen weiter validieren können.
				\end{abstract}
			\end{minipage}
			
			\vspace{0.3cm}
			\noindent
			\small\textbf{Schlüsselwörter:} YUCT, kritische Gruppengröße, Phasenübergang, Bienenschwärmen, militärische Hierarchie, Fischschwarm, universelle Skalierung, $\beta = 0,67$, technologische Schwelle.
		\end{center}
	\end{titlepage}
	
	\tableofcontents
	\newpage
	
	\section{Einleitung}
	\label{sec:intro}
	
	Eine der zentralen Vorhersagen der Yakushev Unified Coordination Theory (YUCT) ist die Existenz eines universellen Exponenten $\beta = 0,67$, der die Skalierung von Koordinationsfehlern mit der Systemgröße bestimmt (Anhang L). Dieser Exponent ergibt sich aus der fraktalen Struktur von Wörterbüchern und informationstheoretischen Beschränkungen und wurde über 40 Größenordnungen hinweg verifiziert – von der DNA-Replikation bis zu Fluktuationen der kosmischen Mikrowellenhintergrundstrahlung. Wenn $\beta$ wirklich universell ist, sollte er sich auch in den kritischen Größen manifestieren, bei denen kollektive Systeme Phasenübergänge durchlaufen, wie Schwärmen, Aufspaltung oder hierarchische Reorganisation. Darüber hinaus besagt das Axiom der technologischen Wörterschwelle (Anhang O), dass ein System ein minimales technologisches Niveau $\Tmin$ erreichen muss, um ein Wörterbuch zu nutzen oder zu erzeugen; diese Schwelle kann als ein kritischer Punkt in der Koordinationseffizienz $\Keff$ des Systems reinterpretiert werden. Dieser Anhang sammelt empirische Daten aus verschiedenen Bereichen und zeigt, dass die beobachteten Verhältnisse von kritischer zu optimaler Gruppengröße mit den Vorhersagen von YUCT übereinstimmen, wodurch das abstrakte Axiom mit konkreten, messbaren Größen verbunden wird.
	
	\section{Theoretischer Rahmen}
	\label{sec:theory}
	
	Für jedes koordinierte System der Größe $N$ wird angenommen, dass der relative Koordinationsfehler $\varepsilon$ skaliert wie
	\begin{equation}
		\varepsilon(N) = \varepsilon_0 N^{\beta}, \qquad \beta = 0,67,
		\label{eq:scaling}
	\end{equation}
	wobei $\varepsilon_0$ eine systemabhängige Konstante ist. Das System arbeitet optimal bei einer Größe $N_{\text{opt}}$, bei der $\varepsilon$ minimal ist. Wenn $N$ über $N_{\text{opt}}$ hinaus wächst, wachsen die Fehler, und wenn sie einen kritischen Schwellenwert $\varepsilon_{\text{crit}}$ erreichen, durchläuft das System einen Phasenübergang (z. B. Aufspaltung in Untergruppen). Somit gilt
	\begin{equation}
		\varepsilon(N_{\text{crit}}) = \varepsilon_{\text{crit}}.
	\end{equation}
	Wenn die Tochtergruppen nach dem Übergang nahe der optimalen Größe liegen, dann folgt aus der Zahlenerhaltung $N_{\text{crit}} \approx k N_{\text{opt}}$ mit $k > 1$. Der genaue Wert von $k$ hängt von der Art des Übergangs ab (symmetrische oder asymmetrische Teilung). Mit (\ref{eq:scaling}) erhalten wir
	\begin{equation}
		k^{\beta} = \frac{\varepsilon_{\text{crit}}}{\varepsilon_{\text{opt}}}, 
		\qquad\text{und damit}\qquad 
		k = \left( \frac{\varepsilon_{\text{crit}}}{\varepsilon_{\text{opt}}} \right)^{1/\beta}.
		\label{eq:k}
	\end{equation}
	Somit wird das Verhältnis $k$ allein durch das Verhältnis von kritischem zu optimalem Fehler bestimmt, potenziert mit $1/\beta$. Wenn $\varepsilon_{\text{crit}}/\varepsilon_{\text{opt}}$ über verschiedene Systeme hinweg ungefähr konstant ist (eine plausible Hypothese, falls die Schwelle durch fundamentale physikalische oder biologische Randbedingungen gesetzt wird), dann sollte auch $k$ universell sein, und sein Wert kann mit empirischen Daten verglichen werden.
	
	\section{Empirische Daten zu kritischen Gruppengrößen}
	\label{sec:data}
	
	Wir haben Daten aus mehreren gut untersuchten Systemen gesammelt, in denen Gruppengrößen und Übergangspunkte dokumentiert sind.
	
	\subsection{Honigbienenvölker (\textit{Apis mellifera})}
	\label{subsec:bees}
	
	Die Imkerei liefert verlässliche Zahlen:
	\begin{itemize}
		\item Optimale Volksgröße (erster Schwarm, „Primärschwarm“): 50\,000–60\,000 Arbeiterinnen (Durchschnittsmasse 7–8 kg, Einzelmasse ca. 0,12 g). \cite{Seeley2010, Winston1987}
		\item Größe vor dem Schwärmen: 80\,000–100\,000 (manchmal mehr). \cite{bees}
		\item Somit $k_{\text{bees}} = N_{\text{crit}}/N_{\text{opt}} \approx 1,5$–$1,67$.
	\end{itemize}
	Das Schwärmen ist asymmetrisch: der Primärschwarm nimmt die alte Königin mit und hinterlässt ein kleineres Restvolk. Das beobachtete $k$ liegt bei etwa 1,6.
	
	\subsection{Militäreinheiten}
	\label{subsec:military}
	
	Die moderne militärische Organisation zeigt eine klare hierarchische Struktur mit wohldefinierten Einheitsgrößen \cite{FM3-90.5}:
	\begin{itemize}
		\item Gruppe/Trupp: 8–12 Soldaten
		\item Zug: 30–50 Soldaten
		\item Kompanie: 100–200 Soldaten
		\item Bataillon: 300–1000 Soldaten
		\item Brigade: 3000–5000 Soldaten
		\item Division: 9000–18000 Soldaten
	\end{itemize}
	Die Verhältnisse zwischen aufeinanderfolgenden Ebenen liegen zwischen 3 und 5, mit einem Schwerpunkt bei 3–4. Für eine Kompanie, wenn sie über etwa 200 hinauswächst, kann sie in zwei Züge aufgeteilt werden. Dies ergibt $k \approx 2$, was kleiner ist als das Verhältnis zwischen den Ebenen. Bessere Daten über tatsächliche Aufspaltungsereignisse wären wünschenswert, aber die Existenz klarer hierarchischer Ebenen selbst deutet auf quantisierte Sprünge in der organisatorischen Komplexität hin.
	
	\subsection{Fischschwärme – Zebrabärbling (\textit{Danio rerio})}
	\label{subsec:fish}
	
	Kontrollierte Experimente zum kollektiven Verhalten von Zebrabärblingen \cite{Tunstrom2013} zeigen einen Phasenübergang von polarisiertem Schwimmen zu ungeordnetem Kreisen mit zunehmender Dichte. Die kritische Dichte entspricht einer Gruppengröße von etwa 50–100 Fischen in einem Standardbecken. Kein genauer $N_{\text{opt}}$ wird angegeben, aber der Übergang tritt auf, wenn die Gruppe eine bestimmte Größe überschreitet. Dies ist konsistent mit der Idee, dass $k$ etwa 2 beträgt.
	
	\subsection{Delfingruppen (\textit{Tursiops} spp.)}
	\label{subsec:dolphins}
	
	Feldstudien an Großen Tümmlern \cite{Connor2000} berichten typische Gruppengrößen von 30–70 Individuen. In Gegenwart von Raubtieren (Haien) können sich die Gruppen zu Herden von mehreren hundert (bis zu 600) zusammenballen. Diese Aggregation ist ein defensiver Übergang, keine Aufspaltung, aber sie veranschaulicht eine kritische Änderung der Koordination. Das Verhältnis der großen Herde zur typischen Gruppe beträgt etwa 5–10, was dem $k$ in der umgekehrten Richtung (Verschmelzung) entsprechen könnte.
	
	\subsection{Nutztierpopulationen – Genetische Erhaltung}
	\label{subsec:livestock}
	
	Daten der Ernährungs- und Landwirtschaftsorganisation (FAO) zu kritischen Populationsgrößen für die genetische Erhaltung \cite{FAO2013}:
	\begin{itemize}
		\item Rinder: kritischer Status bei Population <2000, gefährdet bei 5000–15000, normal >25000.
		\item Schafe: empfohlen 100–250 Mutterschafe für den Genpool-Erhalt.
		\item Schweine: 25 Eber + 100 Sauen.
		\item Hühner: 50 Hähne + 250 Hennen.
	\end{itemize}
	Diese Zahlen spiegeln die minimal lebensfähige Populationsgröße wider, die als kritische Schwelle für die genetische Koordination (Vermeidung von Inzuchtdepression) angesehen werden kann. Interpretiert man $N_{\text{opt}}$ als die Größe für optimale genetische Gesundheit (z. B. 25000 für Rinder) und $N_{\text{crit}}$ als die minimal lebensfähige (z. B. 2000), erhält man $k \approx 0,08$, was kleiner als 1 ist – dies ist eine andere Art von Übergang (Aussterben). Dennoch könnte das Verhältnis in gewisser Weise ebenfalls mit $\beta$ skalieren, was jedoch ein separates Modell erfordert.
	
	\subsection{Primatengruppen}
	\label{subsec:primates}
	
	Daten über Pavian-Gruppen \cite{Alberts1995} zeigen typische Größen von 30–80 Individuen. Wenn eine Gruppe zu groß wird (über etwa 100), kann sie sich in zwei Tochtergruppen aufspalten. Dies ergibt $k \approx 2$–$3$.
	
	\section{Vergleich mit der YUCT-Vorhersage}
	\label{sec:comparison}
	
	Gemäß Gleichung (\ref{eq:k}) hängt das Verhältnis $k = N_{\text{crit}}/N_{\text{opt}}$ vom Fehlerverhältnis $\varepsilon_{\text{crit}}/\varepsilon_{\text{opt}}$ ab. Wenn dieses Fehlerverhältnis universell ist, dann sollte $k$ über verschiedene Arten hinweg konstant sein. Die in Tabelle \ref{tab:summary} zusammengestellten Daten deuten darauf hin, dass $k$ tatsächlich in einem relativ engen Bereich liegt, etwa zwischen 1,5 und 2,5.
	
	\begin{table}[htbp]
		\centering
		\caption{Zusammenfassung der kritischen Größenverhältnisse für verschiedene Systeme}
		\label{tab:summary}
		\begin{tabular}{lccc}
			\toprule
			\textbf{System} & $N_{\text{opt}}$ & $N_{\text{crit}}$ & $k = N_{\text{crit}}/N_{\text{opt}}$ \\
			\midrule
			Honigbienenvolk & 50–60k & 80–100k & 1,5–1,7 \\
			Militäreinheit (Zug→Kompanie) & 30–50 & 100–200 & 2–4 \\
			Zebrabärblingsschwarm & 50? & 100? & 2? \\
			Paviangruppe & 30–80 & ~100 & 1,3–3 \\
			\bottomrule
		\end{tabular}
	\end{table}
	
	Wenn man $k \approx 1,6$ als repräsentativen Wert nimmt (Bienen, Primaten), dann folgt aus Gleichung (\ref{eq:k})
	\[
	\frac{\varepsilon_{\text{crit}}}{\varepsilon_{\text{opt}}} = k^{\beta} = 1,6^{0,67} \approx 1,37.
	\]
	Das bedeutet, dass der kritische Fehler nur etwa 37\% größer ist als der optimale Fehler – ein plausibler Schwellenwert für die Einleitung eines Phasenübergangs. Für $k \approx 2$ (Fische, einige Primatengruppen) erhalten wir $2^{0,67} \approx 1,59$. Die Konsistenz dieser Werte deutet darauf hin, dass $\varepsilon_{\text{crit}}/\varepsilon_{\text{opt}}$ über verschiedene Systeme hinweg tatsächlich ungefähr konstant ist, was die Universalität von $\beta$ und die Interpretation der technologischen/organisatorischen Schwelle als Phasenübergang der Koordinationseffizienz unterstützt.
	
	\section{Überprüfbare Vorhersagen}
	\label{sec:predictions}
	
	Basierend auf dem obigen Rahmen formulieren wir mehrere spezifische Vorhersagen, die experimentell oder beobachtend getestet werden können.
	
	\begin{enumerate}
		\item \textbf{Honigbienen:} Das Verhältnis der Größe des Primärschwarms zur Größe des Volkes vor dem Schwärmen sollte über verschiedene Unterarten und Umweltbedingungen hinweg konstant sein und $k = (\varepsilon_{\text{crit}}/\varepsilon_{\text{opt}})^{1/\beta}$ mit $\beta = 0,67$ erfüllen. Wenn $\varepsilon_{\text{crit}}/\varepsilon_{\text{opt}}$ unabhängig geschätzt werden kann (z. B. aus der Sammel-Effizienz oder Krankheitsinzidenz), sollte das vorhergesagte $k$ mit den Beobachtungen übereinstimmen.
		
		\item \textbf{Fischschwärme:} In kontrollierten Experimenten mit Zebrabärblingen oder anderen Arten sollte die Gruppengröße, bei der der Übergang von polarisiertem Schwimmen zu Kreisen (oder zur Aufspaltung) erfolgt, gemäß demselben Gesetz relativ zur optimalen Gruppengröße skalieren. Messungen von Koordinationsfehlern (z. B. Varianz der Ausrichtung) können unabhängige Schätzungen von $\varepsilon_{\text{opt}}$ und $\varepsilon_{\text{crit}}$ liefern.
		
		\item \textbf{Primaten:} Langzeit-Feldstudien an Pavian-, Makaken- oder Schimpansengruppen sollten Aufspaltungsereignisse sowie die Größen von Eltern- und Tochtergruppen dokumentieren. Das Verhältnis $k$ sollte stabil und mit der YUCT-Vorhersage konsistent sein.
		
		\item \textbf{Militärorganisationen:} Historische Daten über die Aufspaltung von Militäreinheiten (z. B. ein Bataillon teilt sich in zwei) können untersucht werden. Die Größen, bei denen solche Aufspaltungen auftreten, sollten derselben Skalierung folgen, möglicherweise mit einem anderen $k$ aufgrund der hierarchischen Natur.
		
		\item \textbf{Nutztierpopulationen:} Für die genetische Erhaltung könnte die effektive Populationsgröße $N_e$, die zur Aufrechterhaltung der genetischen Vielfalt benötigt wird, mit der kritischen Größe für das Aussterben zusammenhängen. Das Verhältnis der minimal lebensfähigen $N_e$ zur optimalen $N_e$ für das langfristige Überleben könnte ebenfalls mit $\beta$ skalieren, auch wenn dies spekulativer ist.
		
		\item \textbf{Soziale Netzwerke:} Online-Sozialgruppen (z. B. WhatsApp-Chats, Reddit-Communities) teilen sich oft auf, wenn sie eine bestimmte Größe überschreiten. Die kritische Größe für die Aufspaltung und die typische Größe der Tochtergruppen können anhand digitaler Spuren analysiert werden und bieten eine riesige Datenmenge zum Testen der Universalität von $k$ und $\beta$.
	\end{enumerate}
	
	\section{Verbindung zum Axiom der technologischen Wörterschwelle (Anhang O)}
	\label{sec:connection}
	
	Das Axiom der technologischen Wörterschwelle (Anhang O) besagt, dass ein System ein minimales technologisches Niveau $\Tmin$ erreichen muss, um ein Wörterbuch zu nutzen oder zu erzeugen. Im Zusammenhang mit Gruppengrößenübergängen kann $\Tmin$ als die kritische Koordinationseffizienz $K_{\mathrm{eff},\text{crit}}$ reinterpretiert werden, die erforderlich ist, um Kohärenz aufrechtzuerhalten. Das empirische Verhältnis $k = N_{\text{crit}}/N_{\text{opt}}$ wird dann zu einem Maß dafür, wie stark das System wachsen kann, bevor es sich aufspalten oder reorganisieren muss. Die Tatsache, dass $k$ in verschiedenen Systemen in der Größenordnung von 1,5–2 liegt, deutet darauf hin, dass das zugrundeliegende Verhältnis $K_{\mathrm{eff},\text{crit}}/K_{\mathrm{eff},\text{opt}}$ universell ist, im Einklang mit dem fraktalen Skalierungsgesetz. Somit liefern die hier präsentierten Daten eine indirekte Unterstützung für die Existenz einer technologischen Schwelle und deren Verbindung mit dem universellen Exponenten $\beta$.
	
	\section{Fazit}
	\label{sec:conclusion}
	
	Die in diesem Anhang zusammengestellten Daten, wenn auch vorläufig, deuten darauf hin, dass die kritischen Gruppengrößen, bei denen kollektive Systeme Phasenübergänge durchlaufen, nicht willkürlich sind, sondern einem universellen Skalierungsgesetz folgen, das mit dem YUCT-Exponenten $\beta = 0,67$ und dem Axiom der technologischen Wörterschwelle verknüpft ist. Das Verhältnis $k = N_{\text{crit}}/N_{\text{opt}}$ gruppiert sich um 1,5–2,5, und der daraus abgeleitete Fehlerschwellenwert $\varepsilon_{\text{crit}}/\varepsilon_{\text{opt}} \approx 1,3$–1,6 ist über verschiedene Systeme hinweg bemerkenswert konsistent. Wir haben spezifische, überprüfbare Vorhersagen skizziert, die durch zukünftige Experimente und Beobachtungen verifiziert werden können. Die Bestätigung dieser Vorhersagen würde eine starke unabhängige Stütze für YUCT sowie für den fundamentalen Charakter der Koordinationseffizienzmetrik $\Keff$ und ihrer Skalierungsgesetze liefern.
	
	\begin{thebibliography}{99}
		\bibitem{Seeley2010} Seeley, T. D. (2010). \emph{Honeybee Democracy}. Princeton University Press.
		\bibitem{Winston1987} Winston, M. L. (1987). \emph{The Biology of the Honey Bee}. Harvard University Press.
		\bibitem{FM3-90.5} FM 3-90.5 (2008). \emph{Combined Arms Battalion}. US Army Field Manual.
		\bibitem{Tunstrom2013} Tunstrøm, K. et al. (2013). Collective states, multistability and transitional behavior in schooling fish. \emph{PLoS Computational Biology}, 9(2), e1002915.
		\bibitem{Connor2000} Connor, R. C. et al. (2000). The bottlenose dolphin: social relationships in a fission‑fusion society. In: \emph{Cetacean Societies}, University of Chicago Press.
		\bibitem{FAO2013} FAO (2013). \emph{In vivo conservation of animal genetic resources}. FAO Animal Production and Health Guidelines No. 14.
		\bibitem{Alberts1995} Alberts, S. C. \& Altmann, J. (1995). Balancing costs and opportunities: dispersal in male baboons. \emph{American Naturalist}, 145(2), 279–306.
		\bibitem{Yakushev2026} Yakushev, A. V. (2026). \emph{Yakushev Unified Coordination Theory (YUCT)}. Zenodo. \url{https://doi.org/10.5281/zenodo.18444599}
		\bibitem{AppendixL} Anhang L: Fraktale Skalierung des Koordinationsfehlers – Ein universelles Gesetz von der DNA bis zur Kosmologie.
		\bibitem{AppendixO} Anhang O: Das Axiom der technologischen Wörterschwelle.
		% Fehlende Referenz für \cite{bees} hinzugefügt
		\bibitem{bees} Seeley, T. D. (2010) – wie oben, oder durch eine spezifische Referenz ersetzen.
	\end{thebibliography}
	
\end{document}